
\section{讨论：关于“波/粒二分”的严格表达（信息与群轨道语言）}

我们把“不可被其他元素恢复出来的数据算正交”具体化为闭包图在 \WL{}/相干配置意义下的残余轨道（定义 4）。据此，可将经验判断精炼为可检验命题。

\subsection{“粒子态”：塌缩快且残余对称小}

在闭包图上，若存在某类微观邻接使得 $t_{\mathrm{resolve}}$ 在 $m$ 增大时仍保持 $O(1)$ 或 $O(\log\abs{X_m})$，且残余未分解类规模 $s_{\max}$ 受常数上界控制，则该邻接实现了“可控时间的可逆检索”：回溯搜索树被 \WL{} 迅速剪枝，纤维信息可在短深度内被定位。

\subsection{“波态”：塌缩慢或存在大规模未分解}

当 $d\le 2$ 或几何极薄时，\WL{} 信息扩散慢，导致 $t_{\mathrm{resolve}}$ 快速增长甚至在给定步数内无法完全分解；这对应“时间纤维信息在局部统计下长期不可见”。其机制与复返随机游走一致（参见 \cite{novakpolya} 的阈值现象学）。

\subsection{“时间波/概率波”：共振 m 的残余 Z2 时间位}

在少数共振 $m$ 上，即使微观图高对称（如 hypercube），仍可能出现无法由 $1$-\WL{} 消除的 $\ZZ_2$ 方向。它不是噪声，而是闭包图自同构群在折叠相容约束下留下的真实轨道；在本文模型里其可验证对象正是：\WL{} 细化不变的残余配对轨道。

